% ===========================================================================
% Discussion
% File: sections/05_Discussion.tex
% Built sub-section by sub-section
% ===========================================================================

\section{Discussion}
\label{sec:discussion}

% ---------------------------------------------------------------------------

% ---------------------------------------------------------------------------
\subsection{The PAR-Saturation Mechanism: Implications for Site
	Selection and Crop Choice}
\label{sec:discussion_par}
% ---------------------------------------------------------------------------

\hspace{1cm}The non-monotone relationship between annual solar resource
and eLER --- Freiburg (GHI = 1\,150\,kWh\,m$^{-2}$) outperforming
Konya (GHI = 1\,650\,kWh\,m$^{-2}$) --- is the most physically
significant finding of this study and requires careful mechanistic
interpretation before deployment implications can be drawn.

The mechanism operates through the PAR integral denominator in
Equation~\ref{eq:lercrop}. In an open field, the denominator
$\int_{t_0}^{t_1} PAR_{\text{open}}(t)\,dt$ accumulates all incoming
PAR without a ceiling. At Konya, with a mean growing-season daytime
GHI of approximately \SI{450}{W\,m^{-2}}, the corresponding open-field
PAR of $450 \times 0.48 = \SI{216}{W\,m^{-2}}$ exceeds
$PAR_{\text{sat}} = \SI{174}{W\,m^{-2}}$ for the majority of daylight
hours. The denominator therefore grows faster than the numerator ---
which is clipped at $PAR_{\text{sat}}$ in both the APV and the open
field --- making $LER_{\text{crop}}$ structurally lower at
high-irradiance sites regardless of the panel geometry.

At Freiburg, mean growing-season daytime GHI is approximately
\SI{320}{W\,m^{-2}}, giving open-field PAR of
\SI{154}{W\,m^{-2}} --- below $PAR_{\text{sat}}$. The denominator
accumulates proportionally less uncapped PAR, and the ratio remains
more favourable. This is not a statement about absolute crop yield:
Freiburg produces less tomato per hectare in absolute terms
(30\,t\,ha$^{-1}$) than Konya (60\,t\,ha$^{-1}$). It is a statement
about the \textit{relative} efficiency of land use: at Freiburg, the
APV system captures a larger fraction of the agronomically useful
PAR, because the open-field reference does not accumulate as much
PAR above the saturation point.

\textbf{Generalisation to other crops.} The PAR-saturation mechanism
is crop-specific. Tomato has a high light saturation point
($PAR_{\text{sat}} = \SI{174}{W\,m^{-2}}$), making it sensitive to
the mechanism at the irradiance levels of all four study sites. For
shade-tolerant crops with lower saturation points --- lettuce
(\SI{76}{W\,m^{-2}}), spinach (\SI{80}{W\,m^{-2}}), or basil
(\SI{100}{W\,m^{-2}}) --- the denominator would be clipped at a much
lower threshold, and open-field PAR would exceed $PAR_{\text{sat}}$
for a larger fraction of daylight hours at \textit{all} sites,
potentially eliminating the eLER difference between high- and
low-irradiance locations. For high-saturation crops such as maize
($PAR_{\text{sat}} \approx \SI{400}{W\,m^{-2}}$), the mechanism
would be reversed: open-field PAR rarely reaches saturation, the
denominator grows proportionally with irradiance, and high-irradiance
sites would achieve higher $LER_{\text{crop}}$, restoring a monotone
eLER--GHI relationship. Multi-crop generalisation is therefore
identified as the highest-priority extension of the present framework.

\textbf{Site selection implication.} The conventional wisdom in
agrivoltaic deployment --- that high-irradiance semi-arid sites
are preferable because they maximise PV revenue --- requires
qualification when high-PAR crops such as tomato are cultivated.
The present results suggest that moderate-irradiance temperate
sites offer a comparable or superior eLER for tomato, because the
PAR available in those climates is more fully utilisable by the
crop. From a food-energy-water nexus perspective, the optimal site
is not the site with the most sun, but the site where the solar
resource best matches the crop's photosynthetic capacity.
This reframes the site selection problem from a single-objective
(maximise PV output) to a multi-objective (maximise eLER) decision,
consistent with the optimization framework developed in this study.

% ---------------------------------------------------------------------------
\subsection{Boundary Convergence: Physical Meaning and the
	Access Trade-off}
\label{sec:discussion_boundary}
% ---------------------------------------------------------------------------

\hspace{1cm}The convergence of the optimizer to the lower bounds of
both row spacing ($d_{\text{row}}^* = \SI{3.03}{m}$) and module
height ($H_m^* = \SI{1.50}{m}$) at all four sites is a structural
result that reveals the shape of the eLER landscape and has direct
practical implications for APV system design and agricultural
operations.

\textbf{Why the optimizer converges to the boundary.}
Increasing panel density --- reducing $d_{\text{row}}$ and $H_m$
toward their lower bounds --- has two simultaneous effects on eLER.
First, it increases $LER_{\text{PV}}$: more modules per hectare
generates more electricity from the same land area, raising the
numerator $E_{\text{PV,sold}}$ relative to the reference. Second,
it increases $F_{\text{shad}}$, which reduces $PAR_{\text{open}}$
denominator in Equation~\ref{eq:lercrop} by intercepting more
incident radiation. Both effects push eLER upward. The only
countervailing force is the reduction in $LER_{\text{crop}}$: more
shading reduces PAR available to the crop, lowering the numerator
$\int\min(PAR_{\text{AV}}, PAR_{\text{sat}})\,dt$. The optimizer
converges to the physical constraint boundary because, at the
parameter values of this study (tomato, GHI = 1\,150--2\,050
kWh\,m$^{-2}$), the PV density gain and denominator reduction
outweigh the crop yield penalty up to the point where physical
overlap between rows prevents further densification.

\textbf{The eLER penalty of mechanical access.}
In practice, agricultural machinery requires a minimum inter-row
clearance for tractors, harvesters, and irrigation equipment. Typical
tractor widths of 1.8--2.2\,m, combined with operational clearance,
require row spacings of 3.5--5.0\,m in mechanised farming systems.
The eLER cost of this access requirement can be quantified from the
boundary convergence result: since the optimizer seeks the minimum
feasible $d_{\text{row}}$ at all sites, each additional metre of
required row spacing represents a deviation from the eLER optimum.
From the sensitivity of eLER to $d_{\text{row}}$ in the neighbourhood
of the boundary solution, the eLER penalty is approximately
$-$0.04 to $-$0.06 per additional metre of row spacing, with the
exact value depending on site irradiance and tilt angle. A mechanised
farming system requiring $d_{\text{row}} = \SI{5.0}{m}$ instead of
\SI{3.03}{m} would sacrifice approximately 0.08--0.12\,eLER units
--- a penalty of 5--7\% relative to the boundary optimum.

\textbf{Practical design recommendations.}
Three design strategies can partially recover the access penalty
while maintaining agricultural operability:

\begin{enumerate}[leftmargin=*, noitemsep]
	
	\item \textbf{Smallholder manual cultivation.} For systems
	serving smallholder farmers who operate without mechanised
	equipment --- the predominant farming model in Burkina Faso
	and much of the semi-arid developing world --- the minimum
	physical clearance of \SI{3.03}{m} is operationally feasible.
	Manual cultivation, hand harvesting, and drip irrigation
	require only 0.8--1.2\,m of working clearance, making the
	boundary-convergence design directly deployable at Ouagadougou-type
	sites.
	
	\item \textbf{Increased module height.} Raising $H_m$ from
	\SI{1.50}{m} to \SI{2.5}{m} allows agricultural machinery to
	pass beneath the panels even at reduced row spacing, recovering
	mechanical access without increasing $d_{\text{row}}$.
	The eLER cost of this height increase is modest: $H_m$ affects
	$F_{\text{shad}}$ only through the shadow length formula
	(Equation~\ref{eq:shading}), and at $d_{\text{row}} = \SI{3.03}{m}$
	the shading fraction is already near its maximum for the relevant
	solar elevation angles.
	
	\item \textbf{Shade-adapted crop selection.} Choosing crops with
	lower $PAR_{\text{sat}}$ (lettuce, spinach, herbs) reduces the
	eLER penalty of increased row spacing because the crop is less
	sensitive to the absolute PAR reduction under denser panel
	configurations. For such crops, a wider $d_{\text{row}}$ may
	achieve a similar or higher $LER_{\text{crop}}$ as the
	boundary-convergence design for tomato.
	
\end{enumerate}

\textbf{Comparison with prior literature.} The boundary convergence
on row spacing is consistent with the findings of \citet{mazzeo2025},
who reported that LER is monotonically increasing with GCR for
lettuce up to GCR = 0.65 before levelling off, and with
\citet{riaz2022}, who found that the LPF-optimal design for bifacial
tomato in Arizona converged to the maximum permissible GCR under
their access constraints. The present study extends these findings
by providing an explicit quantification of the access penalty and
by demonstrating that the boundary convergence is climate-independent
--- it occurs at all four sites despite GHI varying by a factor of
1.8 and mean annual temperature varying by 17\,\si{\celsius}.

% ---------------------------------------------------------------------------
\subsection{The eLER Metric: Interpretation, Definitions,
	and Structural Properties}
\label{sec:discussion_eler}
% ---------------------------------------------------------------------------

\hspace{1cm}Three structural properties of the eLER metric require
explicit interpretation to avoid misreading the numerical results
reported in Section~\ref{sec:results_eler}.

\textbf{Property 1 --- $LER_{\text{biogas}} < 1$ is a coupling
	consequence, not a system failure.}
The result $LER_{\text{biogas}} = 0.615$--$0.727$ at all sites
means that the APV--AD system produces less biogas than a standalone
reference AD unit processing the same land area under open-field
conditions. This is not a design failure or an inefficiency of the
integrated system: it is a direct and unavoidable consequence of the
shading-residue coupling. When panels shade the crop, $LER_{\text{crop}}
< 1$, meaning fewer crop residues are harvested per hectare. Fewer
residues reduce VS input to the digester, which reduces biogas output.
The reference AD unit, by contrast, operates with open-field residues
at $LER_{\text{crop}} = 1.0$ --- maximum residue production.
The ratio of the two biogas outputs is therefore structurally less
than 1.0 whenever shading is present.

This coupling has an important implication for the interpretation of
the eLER sum. The three eLER components are not independent:
$LER_{\text{crop}}$ and $LER_{\text{biogas}}$ are negatively
correlated through the shading-residue pathway. Increasing GCR
raises $LER_{\text{PV}}$ and reduces $LER_{\text{crop}}$, which in
turn reduces $LER_{\text{biogas}}$. The optimizer finds the design
that maximises the sum of all three components, accepting the
biogas penalty in exchange for the larger PV gain.

\textbf{Property 2 --- An alternative $LER_{\text{biogas}}$
	definition changes the interpretation but not the system physics.}
An alternative reference for $LER_{\text{biogas}}$ could use a
standalone AD unit operating \textit{without} PV heating --- at
ambient temperature rather than at the PV-assisted
$T_{\text{dig,eff}}$. Under this definition:
\begin{equation}
	LER_{\text{biogas,alt}} =
	\frac{E_{\text{biogas,APV}}}{E_{\text{biogas,ambient}}}
	\label{eq:lerbiogas_alt}
\end{equation}
where $E_{\text{biogas,ambient}}$ is computed using
$T_{\text{dig,eff}} = T_{\text{amb}}$ without the 5\,\si{\celsius}
boost or the $[26,\,37]$\,\si{\celsius} clip. Under this definition,
$LER_{\text{biogas,alt}} > 1$ at Konya and Freiburg in winter months,
explicitly quantifying the PV thermal subsidy to AD performance. This
alternative is conceptually more informative about the PV--AD synergy
but harder to compare across studies because the reference changes
with ambient temperature. The primary definition used throughout this
paper (Equation~\ref{eq:lerbiogas}) was chosen for consistency with
the classical LER framework, in which all references are
activity-specific open-field monocultures operating under the same
boundary conditions as the integrated system.

\textbf{Property 3 --- ESR $\gg$ 1 is a design identity, not a
	performance finding.}
ESR values of 533--1\,413 at the optimised designs reflect the
fundamental size mismatch between the PV array (which generates
electricity at the scale of MWh\,ha$^{-1}$\,yr$^{-1}$) and the
small farm-scale digester (which demands heat at the scale of
MWh\,ha$^{-1}$\,yr$^{-1}$ --- the same unit but three orders of
magnitude smaller in absolute value). With $V_{\text{dig}}^* =
3.10$--$4.14$\,m$^3$ and heat demand of only
1.12--2.08\,MWh\,ha$^{-1}$, even 5\% of PV output provides
55--79\,MWh of thermal energy --- 27--70 times the annual heat
demand. The ESR is therefore more a confirmation of the PV array
being appropriately scaled to the land area than a meaningful
performance indicator for the AD system. Future studies operating
at larger scale (community digesters serving multiple hectares)
would achieve lower ESR values that are more informative about the
thermal balance between PV and AD.

\textbf{Definition A versus Definition B: when to use each.}
Definition~A ($LER_{\text{PV}}$ referenced to a dense-pack
independent array) is the scientifically rigorous choice for
studies that aim to quantify the net land-use benefit of the APV
configuration relative to what could be achieved with two separate
monocultures. Definition~B ($LER_{\text{PV}} \equiv GCR$) is
appropriate only for literature comparison with studies that have
used the circular reference, and its numerical values should always
be flagged as such. The difference between the two definitions is
0.031--0.032 in $LER_{\text{PV}}$ and 0.031--0.032 in eLER at the
optimised designs of this study (Table~\ref{tab:eLER}), which is
small relative to the between-site eLER variation (0.105) but
non-negligible relative to the between-site $LER_{\text{crop}}$
variation (0.034). Standardisation of the PV reference definition
--- ideally through an agreed protocol analogous to the IEC 61724
standard for PV performance reporting --- is identified as a
priority for the agrivoltaic community.

% ---------------------------------------------------------------------------
\subsection{Economic Viability and Policy Implications}
\label{sec:discussion_economics}
% ---------------------------------------------------------------------------

\hspace{1cm}The 4E analysis reveals a wide range of economic
performance across the four sites, from strongly viable
(Ouagadougou, IRR = 31.1\%) to marginal but positive (Konya,
IRR = 11.8\% against an 8\% discount rate). Three policy-relevant
implications emerge from the economic results.

\textbf{Electricity tariff structure is the primary economic
	determinant.} The dominant revenue stream at all four sites is PV
electricity sales, which accounts for 73--84\% of gross revenue
under Scenario S0. The electricity price therefore has a
disproportionate influence on project economics relative to tomato
price, biogas price, or carbon credit. At Konya, the low Turkish
electricity tariff (\$0.06\,kWh$^{-1}$) is the binding constraint
on IRR: with the same PV output as Almer\'{i}a (which also receives
approximately 1\,500\,MWh\,ha$^{-1}$\,yr$^{-1}$ at optimal design),
Konya generates 33\% less PV revenue simply because the tariff is
33\% lower (\$0.06 vs.\ \$0.09\,kWh$^{-1}$). A sensitivity
analysis shows that a Turkish electricity price increase to
\$0.09\,kWh$^{-1}$ --- consistent with the regional trend as
fossil fuel subsidies are progressively reduced --- would raise
Konya's IRR to approximately 16\%, comparable to Freiburg and
well above the discount rate. This underscores that APV--AD
investment viability in semi-arid developing economies is primarily
a function of energy price reform, not of solar resource or
technology cost.

\textbf{The carbon credit scenario (S6) is the dominant economic
	strategy and should be prioritised in policy design.} S6 raises
IRR by 1.7--5.0 percentage points above S0 at all sites, with the
largest absolute improvement at Freiburg (+5.0\,pp) driven by the
EU ETS price (\$65\,tCO$_2^{-1}$). At Ouagadougou and Konya,
where voluntary carbon market prices are much lower
(\$12--15\,tCO$_2^{-1}$), the S6 improvement is smaller in
absolute terms (+2.0 and +1.7\,pp respectively) but still
meaningful relative to the discount rate. The policy implication
is that extending carbon credit mechanisms --- whether through
the Article 6 Paris Agreement framework, Gold Standard
certification, or national carbon pricing --- to smallholder
agrivoltaic installations in Sub-Saharan Africa and Central
Asia would provide a material improvement in project economics
without requiring technology subsidies or feed-in tariff reform.
Given that the 951\,tCO$_2$\,ha$^{-1}$\,yr$^{-1}$ avoided at
Ouagadougou represents a significant fraction of the annual per
capita CO$_2$ budget compatible with 1.5\,\si{\celsius} pathways,
this mechanism also aligns economic incentives with climate goals.

\textbf{Ouagadougou provides the strongest investment case and the
	most compelling argument for APV--AD deployment in Sub-Saharan Africa.}
An IRR of 31.1\% under S0 and 33.1\% under S6, combined with a
payback of 3.0--3.2\,yr, makes the Ouagadougou configuration
attractive even at the 10\% discount rate that reflects the higher
cost of capital in developing economies. Three structural advantages
drive this result: the highest electricity price in the study
(\$0.12\,kWh$^{-1}$, reflecting the true cost of diesel generation),
the highest biogas yield (5.99\,MWh\,ha$^{-1}$, from the optimal
Arrhenius BMP at 33.2\,\si{\celsius} average digester temperature),
and the largest CO$_2$ avoidance (951\,tCO$_2$\,ha$^{-1}$, from the
high grid emission factor of 0.632\,tCO$_2$\,MWh$^{-1}$). The small
digester volume (3.52\,m$^3$) and low AD CAPEX (\$2\,816) mean that
the biogas revenue stream is essentially free cash flow on top of
the PV investment. For development finance institutions seeking
bankable climate-positive agricultural projects in the Sahel,
the APV--AD model at Ouagadougou-type sites offers a rare
combination of high IRR, short payback, large CO$_2$ impact,
and food security co-benefit.

\textbf{LCOE is competitive with utility-scale solar at the best
	sites.} LCOE$_{\text{PV}}$ of \$0.034\,kWh$^{-1}$ at Almer\'{i}a
approaches the IRENA 2023 global utility-scale solar benchmark of
\$0.049\,kWh$^{-1}$, and the \$0.043--0.045\,kWh$^{-1}$ range
at the other three sites remains within 10--20\% of this benchmark.
The premium relative to utility-scale ground-mount reflects two
factors: the elevated APV structure (\$0.63\,W$_p^{-1}$ vs.\
\$0.45\,W$_p^{-1}$ for ground-mount~\citep{irena2022}) and the
inter-row shading loss at GCR = 69--72\%, which reduces annual
specific yield below what a dense-pack unshaded array would
achieve. However, this LCOE premium purchases the additional
outputs of crop production, water savings, and biogas --- none
of which are captured in the single-output LCOE metric. A
full value-stack LCOE that credits all outputs would be lower
than the PV-only LCOE at all four sites, and likely competitive
with or below the utility-scale benchmark at Ouagadougou and
Almer\'{i}a.

% ---------------------------------------------------------------------------
\subsection{Limitations and Future Work}
\label{sec:discussion_limits}
% ---------------------------------------------------------------------------

\hspace{1cm}Six limitations of the present modelling approach are
identified, each motivating specific future research directions.

\textbf{L1 --- BMP temperature model accuracy (V2 MAPE = 26.9\%).}
The single-exponential Arrhenius model systematically underestimates
BMP by 30--52\% at temperatures of 20--30\,\si{\celsius}, the range
relevant to cold-season operation at Konya and Freiburg. The root
cause is that a single Arrhenius coefficient calibrated at
37\,\si{\celsius} cannot capture the inflected microbial response
curve in the mesophilic lower range, where community adaptation
provides higher-than-exponential yields~\citep{zwietering1990}.
Future work should implement a piecewise Arrhenius model with
separate coefficients above and below 30\,\si{\celsius}, or a
modified Gompertz kinetic model, calibrated against the
multi-temperature BMP datasets now available for
tomato residues and cattle manure~\citep{feng2013, pilarski2025}.
Until such a model is validated, all biogas yield and economic
results for cold-climate sites should be treated as conservative
lower bounds, with the true values likely 15--30\% higher.

\textbf{L2 --- Single-crop scope.}
All results are conditioned on tomato (\textit{Solanum lycopersicum})
with $PAR_{\text{sat}} = \SI{174}{W\,m^{-2}}$. The PAR-saturation
mechanism that drives the non-monotone eLER--GHI gradient is
crop-specific, and the magnitude and direction of this effect will
differ substantially for other crops. Shade-tolerant crops with
lower $PAR_{\text{sat}}$ (lettuce at \SI{76}{W\,m^{-2}}; spinach
at \SI{80}{W\,m^{-2}}) would likely eliminate or reverse the
Freiburg--Konya eLER inversion, while high-saturation crops such
as maize (\SI{400}{W\,m^{-2}}) would restore a monotone eLER--GHI
relationship. A multi-crop sensitivity analysis --- requiring only
a change of the $PAR_{\text{sat}}$ parameter in the existing model
--- is identified as the highest-priority extension, as it would
transform the PAR-saturation finding from a tomato-specific result
into a general mechanistic framework applicable to any crop species.

\textbf{L3 --- Evapotranspiration model simplification.}
Reference ET$_0$ is estimated using the Hargreaves--Samani equation,
which agrees with FAO-56 Penman-Monteith to within $\pm$10\% in
semi-arid and temperate climates~\citep{droogers2002} but cannot
capture humidity and wind effects on ET under the APV canopy.
The ET reduction coefficient $\alpha_{\text{shade}} = 0.40$ is
taken from field measurements under similar APV configurations
\citep{elamri2018, tawassulli2025} but may vary by $\pm$0.10--0.15
depending on crop canopy closure, soil type, and wind exposure.
A $\pm$25\% variation in $\alpha_{\text{shade}}$ propagates to a
$\pm$25\% variation in $W_{\text{saved}}$ and $LER_{\text{water}}$,
which does not affect eLER or economics but is relevant to
water-stressed deployment contexts. Future work should implement
FAO-56 Penman-Monteith with auxiliary humidity and aerodynamic
resistance terms derived from the PVGIS wind speed column, and
validate $\alpha_{\text{shade}}$ against lysimeter measurements
at each study site.

\textbf{L4 --- Absence of digestate-to-crop nutrient feedback (G6).}
The model treats the AD digester as a terminal sink for crop
residues: digestate is valued only as a tradeable nitrogen credit
(Scenario S4) without modelling its effect on soil nutrient
availability and subsequent crop yield. In practice, returning
digestate to the APV field closes a nutrient loop that can
increase fruit yield by 10--25\% relative to unfertilised
conditions~\citep{lallement2023}, which would raise both the
crop revenue and the VS input to the digester in subsequent
seasons. Modelling this feedback requires multi-season soil
nutrient dynamics that are beyond the scope of the present
single-year TMY-based framework. It is identified as Gap G6
and the highest-priority direction for experimental field
validation, requiring co-located APV--AD installations with
controlled digestate application trials over at least three
growing seasons.

\textbf{L5 --- Static economic parameters.}
All economic results use fixed prices for electricity, tomato,
biogas, and carbon credits, without modelling price trajectories
over the 20-year project life. PV module costs have declined at
approximately 10--15\% per year over the past decade~\citep{irena2022},
and biomethane prices are subject to regulatory risk in all four
jurisdictions. A stochastic techno-economic model incorporating
Monte Carlo sampling of price distributions --- including correlated
electricity and carbon price scenarios consistent with IEA
Net Zero pathways~\citep{iea2025} --- would provide more realistic
NPV confidence intervals and identify the price combinations
under which each site crosses the investment threshold.
Such a model would be particularly valuable for Konya, where
the S0 IRR margin above the discount rate is narrow (3.8\,pp)
and highly sensitive to electricity price assumptions.

\textbf{L6 --- Single-hectare normalisation and scalability.}
All results are normalised to 1\,ha, which is appropriate for
smallholder farming contexts (Ouagadougou, Konya) but may not
represent the economics of larger commercial APV--AD installations.
At larger scale (10--100\,ha), several cost components change
non-linearly: AD digester CAPEX per m$^3$ decreases with scale
(economies of scale in tank construction and gas handling),
BOP costs per kWp decrease with larger inverter units, and the
ESR --- currently 533--1\,413 --- would decrease toward more
physically meaningful values as the digester is sized to process
residues from multiple hectares. A scale-up analysis from
smallholder (1\,ha) to commercial (50\,ha) would provide the
full range of deployment contexts and is a natural next step
following experimental validation of the single-hectare model.